\documentclass{article}
\usepackage{amsmath, amssymb}
\setlength{\parindent}{0pt}

\begin{document}

\textbf{CS 2050 --- Notes 09/16}

\bigskip

\textbf{Vacuous Proof}

\medskip

If we have $p\Rightarrow q$ and can prove that $p$ is false, then $p\Rightarrow q$ is automatically true. This is a vacuous proof.

\bigskip

\textbf{Trivial Proof}

\medskip

Whenever $q$ is true, $p\Rightarrow q$ is true.

\bigskip

\textbf{Proof of Equivalence}

\medskip

To prove that a theorem is a biconditional, prove that both $p\Rightarrow q$ and $q\Rightarrow p$ are true.

\end{document}
